\section{Realizing a test on its original Galois lattices}

% BEGIN R85P1
\paragraph{R85P1. Coefficients, Frobenius and maps on the same crystal.}
Use the common tame field $K_t$ and the enlarged coefficient integers
$O$ from Q84T1--2. Write $W=W(\mathbb F_{p^a})$, $A_0=W[[u]]$ and
$E=u^e+p$. Take a good lattice $T$ restricted to $G_{K_t}$.
We use the crystalline-lattice equivalence in Bhatt--Scholze,
\emph{Prismatic $F$-crystals and crystalline Galois representations},
Theorems 5.6 and 7.9. It gives a crystal $\mathscr E$ on $O_{K_t}$,
its integral evaluation $M$ on $(A_0,(E))$, and the original
\'{e}tale lattice $T$. The endomorphisms giving $T$ its
$O$-action lift uniquely as maps of crystals. The ring
$W\otimes_{\mathbb Z_p} O$ splits over the $a$ unramified embeddings.
Each component $M_\tau$ is finite over $A=O[[u]]$ and free there.
Indeed $M$ is finite free over the local $A_0$ and $p,u$ form a
regular sequence there, hence also on each direct summand. Write
$p=\lambda^h$ times a unit in $O$. Multiplication by $\lambda^{h-1}$
embeds $M_\tau/\lambda$ in $M_\tau/p$. Thus $\lambda,u$ is a
regular sequence on $M_\tau$; over the regular local ring $A$ of
dimension two, the depth formula gives projective dimension zero
for each nonzero component. The zero components are free of rank zero.
The Frobenius is an isomorphism after inverting $E$, permutes
the components cyclically and gives equality of their ranks.
Equivalently these ranks are the $O$-rank of $T$ by evaluation in
$O\otimes W(C_p^\flat)$, with its corresponding idempotents.
An $O$-linear equivariant map of crystalline lattices lifts uniquely,
respects the component decompositions, and commutes with Frobenius
and with every crystal base change. This works simultaneously on
the kernel, middle, target and the crystalline extensions used below.
On the integral Witt evaluation the maps are scalar extensions of
the actual maps of $T$. Q84T3 therefore gives middle exactness and
a finite terminal cokernel even though the original lattice map is onto.
% END R85P1

% BEGIN R85P2
\paragraph{R85P2. Descent and the common period comparison.}
The tower $\pi_{j+1}^p=\pi_j$, $\pi_0^e=-p$, gives a map of
the Breuil--Kisin prism to $A_{\inf}$ by $u\mapsto[\pi^\flat]$.
An element of the finite tame group consists of a residue-field
automorphism and a multiplier $\omega\in\mu_e$ on $\pi_0$.
Extend it to the tower by multiplying $\pi_j$ by the unique
$p^j$-th root of $\omega$ in $\mu_e$. These substitutions respect
all defining relations (the polynomials at each level are
Eisenstein over the unramified base) and give a group action.
Extend any such automorphism to the algebraic closure. On the
prism the lift acts on $W$ by residue Frobenius to the indicated
power and sends $u$ to the Teichm\"uller multiple dictated by the
tower. This commutes with Frobenius and sends $E$ to $E$.
Its action on $A_{\inf}$ commutes with tower evaluation. Pull back
$\mathscr E$ along the base automorphism of $K_t$. The action
on the original $G_{\mathbb Q_p}$-lattice identifies the two
\'{e}tale realizations with their conjugate $G_{K_t}$-actions.
Full faithfulness gives the corresponding isomorphism of crystals.
At the tower prism this is the desired semilinear action on $M$.
Two extensions differ by an element fixing the entire tower; it
acts trivially on the image of $M$ because it fixes the prism map
and the natural comparison takes the diagonal Witt action to the
crystal action. Faithfulness of integral Witt evaluation gives the
group law, the specified descent on every vector in $M[1/E]$ and
compatibility with morphisms. For $g\in G_{K_t}$, crystal
functoriality on $A_{\inf}$ preserves its integral module and
therefore its integral Frobenius image. Combining with a tower
lift for its residue class proves preservation of that image for
every $g\in G_{\mathbb Q_p}$.

Here is how the comparison uses the very same Witt realization.
In Proposition 5.3 of Bhatt--Scholze the identification of the
crystal with $T$ over $W(C_p^\flat)$ restricts to an identification
after inverting the cyclotomic element on $A_{\inf}$ and extends
to $B_{\rm cris}$. Its target on the crystalline side is the
crystal at $W$ tensored with $B_{\rm cris}$, and the comparison
is $G_{K_t}$-equivariant and Frobenius compatible. The paper uses
the Frobenius translate of the untilt map in its prism convention.
Transport the perfect prism by the Witt Frobenius, so the structural
map in our coordinates is $\theta$ and the Breuil--Kisin generator
specializes to $[\pi^\flat]$. The element $[\epsilon]-1$ and its
Frobenius inverse both become units in $B_{\rm cris}$. On the
Frobenius-column basis the comparison is precomposed with the actual
Frobenius matrix, which is invertible there. This gives the map $j_B$
of the realized datum. In particular any equality with a common
$A_{\inf}$-multiple in the Witt coordinates evaluates to the same
equality in $B_{\rm cris}$, by the restricted comparison before
scalar extension. The horizontal section $\xi$ specializes at
$u=0$ to the identity. It is the Frobenius-compatible trivialization
of the Frobenius columns, with the constant Frobenius at zero used
to identify the two fibres. Uniqueness follows by iteration, since
Frobenius multiplies the $u$-adic order of a difference by $p$.
Equivalently this is the unique trivialization in the proof of
Theorem 7.9 of Bhatt--Scholze. More explicitly, the map of E85C1 from the tower prism to
$(S_W,(p))$ uses the coefficient Frobenius and $u\mapsto u^p$.
It is followed by evaluation of the divided-power variable at
$[\pi^\flat]$ in $A_{\rm cris}$: its $e$-th power is
$E([\pi^\flat])-p$ in the divided-power ideal with $p$ adjoined.
The factorial expansions converge there by completeness. Modulo
$p$ the structural uniformizer $u^p$ has vanishing $e$-th power.
A sufficiently high Frobenius therefore factors its structural
map through the residue field. Pullback along this factorization
and inversion of crystalline Frobenius after inverting $p$ give
the trivialization by the special fibre as in Proposition 5.3.
It has the prescribed Frobenius and value at zero. Uniqueness
identifies it with the section just constructed before specialization
to $A_{\rm cris}$. Thus $j_B$ takes each such section
to the corresponding crystalline constant, which is fixed by
$G_{K_t}$. All identifications commute with the lattice maps.
% END R85P2

% BEGIN R85P3
\paragraph{R85P3. The physical lattice at the Eisenstein divisor.}
Put $D=D_{{\rm cris},K_t}(T[1/p])$. In horizontal coordinates,
the Frobenius columns span the unmodified lattice at the divisor
$E=0$. The physical lattice is its filtered modification. More
explicitly, in the completion at any root of $E$, split the
de Rham filtration of $D\otimes K_t$ and let $d_j$ denote a
basis vector of filtration degree $j$. The completed physical
lattice is the span of $E^{-j}d_j$. Here coefficients in the
residue field of the root have their constant lifts: over a
splitting field the root is a constant and the local parameter
can be taken to be $u-\pi$ or $E$. This description follows from
Construction 6.5 and the recovery of the filtered module in the
proof of Theorem 7.9 of Bhatt--Scholze. The modification is
$\mathrm{Fil}^0(D_{K_t}\otimes B_{\rm dR})$; its tensor-product
filtration gives exactly the displayed generators. The Frobenius
twist of the unramified constants only permutes our indices.
In the unmodified horizontal frame the Euler connection is
$u\,d/du$. Therefore $E\nabla$ preserves the modified lattice:
on $E^{-j}d_j$ it acts by the regular multiplier $-juE'$.
The entries of this operator on integral physical basis vectors
are analytic on the $E$-disc away from its roots and have only
finite poles at them. Indeed the coordinate change in Frobenius
columns is analytic invertible on a larger disc by the normalized
product in K85F4; the remaining Frobenius matrix and inverse
have only $E$-power denominators. The completed calculation
removes the poles, giving $E\nabla M\subset M_{\rm an}$.

For filtration degree $i$, reduction of $\mathcal L\cap E^iM$
in rational analytic coordinates modulo $E\mathcal L$ picks out
the basis elements with $j\ge i$. This can be checked at all
roots simultaneously: reduction modulo any fixed power of $E$
is a finite collection of jets after splitting, and interpolation
by a polynomial gives arbitrary jets in those fibres. One can
also use the analytic filtration test in K85F3, whose division
argument uses only finite upper and lower powers of $E$ for
the Frobenius matrix and inverse. Such powers already exist
before their optimal exponents are known. Hence the jump set is
contained in the prescribed $U$ by Q84T2 and de Rham base change.
The same basis calculation gives $E^b\mathcal L\subset M$ and
$E^HM\subset\mathcal L$ at the completed rational divisor.
These inclusions hold integrally: the entries to test are in
$A[1/E]$, and if their numerators are divisible by the requisite
powers at every $E$-root then the monic remainders vanish over
$O$. For inertia, R85P2 intertwines reduction at zero with
crystalline periods. A nonzero vector of character $c$ in a
component gives a nonzero crystalline period with the same
generator eigenvalue. The period annihilator
$\prod_{s\in\Sigma}(X-\omega^s)$ forces $c\in\Sigma$
because the generator is primitive of order $e$.

Finally the ordinary derivative on the divided-power Frobenius
columns is integral by the first-order crystal calculation E85C1.
Together these properties give the specified heights, characters,
jumps, descent, stability and horizontal periods for each lattice
of R85P1. The crystal functor identifies the unit lattice with the
rank-one unit module. This constructs the realizations needed in
Q84T3--4 from the usual crystalline and Fontaine period inputs.
% END R85P3
